\section{Wnioski}

\begin{itemize}
    \item \textbf{Złożoność obliczeniowa algorytmów}: zaimplementowane algorytmy prezentują różne rzędy złożoności
        obliczeniowej. Najbardziej obciążające są algorytmy dotyczące cykli Hamiltona, które mają złożoność rzędu silni. Takie algorytmy mogą sobie nie radzić z dużymi grafami, dlatego zaleca się stosowanie algorytmów aproksymacyjnych. Sprowadzając algorytmy do złożoności wielomianowej, można analizować duże grafy w dość krótkim czasie, nie tracąc dużo na przydatności i dokładności wyniku.
    \item \textbf{Zastosowanie w praktyce}: wiele problemów praktycznych można sprowadzić do postaci analizy grafu. Korzystając z zaimplementowanych algorytmów, możemy rozwiązywać problemy optymalizacyjne oraz analizować struktury grafów.
\end{itemize}
